\chapter{Conclusions and Future Work}
\label{ch:conclusion}

\section{Summary of Contributions}

This thesis presented a hybrid physics-transformer surrogate modeling framework for train system dynamics. The key contributions are:

\begin{enumerate}
    \item A modular architecture combining reduced-order physics models with transformer-based neural networks, enabling efficient probabilistic forecasting while maintaining physical plausibility
    \item A probabilistic output formulation supporting quantile regression and exceedance probability estimation for safety-critical applications
    \item Comprehensive validation demonstrating improved accuracy and calibration compared to physics-only baselines across diverse scenarios
    \item Ablation studies quantifying the contributions of physics integration, attention factorization, and calibration losses
\end{enumerate}

The framework achieves orders-of-magnitude speedup over high-fidelity simulation while providing calibrated uncertainty estimates essential for risk-aware decision making in railway operations.

% TODO: Refine summary based on final results
% TODO: Add quantitative summary of improvements

\section{Future Work}

\subsection{Real-World Calibration and Validation}

A critical next step is validating the framework on operational data from real train systems. This includes:
\begin{itemize}
    \item Collecting field data on energy consumption, arrival times, and coupler forces
    \item Assessing model accuracy and calibration on real-world scenarios
    \item Identifying and addressing domain gaps between simulation and reality
\end{itemize}

% TODO: Discuss data collection challenges
% TODO: Propose validation methodology

\subsection{Distributed Power Systems}

Current work focuses on single-locomotive consists. Extending to distributed power (multiple locomotives at different positions) requires:
\begin{itemize}
    \item Modeling coordination between locomotives
    \item Handling variable consist topology (locomotives can be anywhere)
    \item Accounting for communication delays in distributed control
\end{itemize}

% TODO: Reference distributed power literature
% TODO: Discuss architectural modifications needed

\subsection{Tighter Safety Constraints}

Integrating regulatory safety requirements directly into the model could enable:
\begin{itemize}
    \item Hard constraints on coupler forces (never exceed thresholds)
    \item Speed limit compliance guarantees
    \item Integration with positive train control (PTC) systems
\end{itemize}

This may require constrained optimization formulations or safe learning approaches.

% TODO: Review safe learning literature
% TODO: Propose constraint integration methods

\subsection{Integration with Optimization and Model Predictive Control}

The surrogate model could be embedded in optimization loops for:
\begin{itemize}
    \item Optimal control plan generation (minimize energy subject to safety constraints)
    \item Real-time model predictive control (MPC) for train operation
    \item Scenario-based optimization under uncertainty
\end{itemize}

This requires differentiable models and efficient gradient computation.

% TODO: Discuss differentiable physics-ML hybrids
% TODO: Reference MPC literature for train systems

\subsection{Productization Considerations}

While this thesis focuses on the internal modeling framework and validation methodology, productization represents a parallel engineering outcome. Key considerations include:
\begin{itemize}
    \item \textbf{Software engineering}: Robust APIs, error handling, logging
    \item \textbf{Deployment}: Integration with existing rail lab infrastructure
    \item \textbf{Maintenance}: Model versioning, monitoring, retraining pipelines
    \item \textbf{User experience}: Intuitive interfaces for scenario definition and result visualization
\end{itemize}

These aspects are important for practical adoption but are outside the scope of this research-focused thesis.

% TODO: Keep this section brief as requested (1 short section only)

\section{Concluding Remarks}

The hybrid physics-transformer framework demonstrates the potential of combining physical knowledge with modern deep learning for efficient, accurate, and uncertainty-aware train dynamics modeling. While challenges remain in generalization, calibration, and real-world validation, the approach provides a foundation for next-generation decision support tools in railway operations.

The integration of physics and ML, probabilistic outputs, and attention-based sequence modeling offers a path toward faster-than-real-time simulation with reliable uncertainty quantification—capabilities essential for both safety analysis and operational optimization in the rail industry.

% TODO: Refine concluding remarks based on final results
